%% kapa: exam=1 round=19 year=2008 period=2 form=A subject=economics subject_ko=경제원론 seq=18 number=18
%% source: Q-Net 감정평가사 기출문제 (cst00309) artlSeq=1001962 file=extracted/1008/1001962/2008년 제19회 감정평가사 1차시험 기출문제/2교시.hwp
%% conversion: kordoc 4.24.51 -> tools/gen1.py (auto); review=proofread-sonnet; flags=none
\begin{question}{18}
다음 중 리카도의 대등정리(Ricardian Equivalence Theorem)가 성립하는 경우로 가장 타당한 것은?
\begin{choices}
  \item 소비자들이 합리적 기대를 하는 경우 성립한다.
  \item 근시안적인 소비자가 존재할 때 성립한다.
  \item 소비자가 현재의 예산제약만을 고려할 때 성립한다.
  \item 유동성제약에 걸린 소비자들이 존재할 때 성립한다.
  \item 정부가 현재의 예산제약만을 고려할 때 성립한다.
\end{choices}
\end{question}
\answer{1}
